\documentclass[12pt,a4paper]{article}
\newcommand{\StudentName}{Vu Tien Giang}
\newcommand{\StudentID}{2570188}
\newcommand{\ReportVariant}{Phasor method and engineering checks}
\usepackage{ee5430_style}

\begin{document}
\MakeTitlePage
\tableofcontents
\newpage

\section{Preliminaries}

All phasors in this report use $\ee^{\jmath\omega t}$.  Thus a forward wave
travelling in the $+z$ direction contains $\ee^{-\gamma z}$.  For a uniform
line,
\begin{equation}
Z_0=\sqrt{\frac{R+\jmath\omega L}{G+\jmath\omega C}},
\qquad
\gamma=\alpha+\jmath\beta
=\sqrt{(R+\jmath\omega L)(G+\jmath\omega C)}.
\end{equation}
The principal square roots are selected so that $\alpha>0$ and
$\Re\{Z_0\}>0$.  Peak (not RMS) amplitudes are used in the time-domain
expressions.

\section{Problem 1}

\subsection{Complex line parameters}

\LineElementDiagram

Convert all data to SI units:
\[
(R,L,G,C)=(1.5,\,10\times10^{-9},\,0.4\times10^{-3},\,
0.2\times10^{-12}),qquad f=2\times10^9~\mathrm{Hz}.
\]
At $\omega=1.256637061\times10^{10}~\mathrm{rad/s}$, define
\[
Z'=R+\jmath\omega L=1.5+\jmath125.663706,
\qquad
Y'=G+\jmath\omega C=0.0004+\jmath0.002513274.
\]
Then
\begin{align*}
\frac{Z'}{Y'}&=4.884148\times10^4+\jmath7.179858\times10^3,\\
Z'Y'&=-0.31522734+\jmath0.05326548.
\end{align*}
Taking the principal square roots,
\begin{align}
Z_0&=221.629712+\jmath16.196087~\Omega,\\
\gamma&=0.047946678+\jmath0.563494654~\mathrm{m}^{-1}.
\end{align}
In polar form, $Z_0=222.220706\angle4.179588\degree~\Omega$.  The requested
wave quantities follow directly:
\begin{align}
\alpha&=0.047946678~\mathrm{Np/m}
=0.41646485~\mathrm{dB/m},\\
\beta&=0.563494654~\mathrm{rad/m},\\
\lambda&=\frac{2\pi}{\beta}=11.1503903~\mathrm{m}.
\end{align}

\checkmarktext\ Multiplying the phasor by $\ee^{-\alpha z}$ shows that the
amplitude decreases by $\ee^{-0.04795}=0.95318$ over one metre, equivalent to
$20\log_{10}(1/0.95318)=0.4165$ dB.

\answer{$Z_0=222.22\angle4.18\degree~\Omega$ and
$\gamma=(0.04795+\jmath0.56349)~\mathrm{m}^{-1}$.  Hence
$\lambda=11.1504~\mathrm{m}$.}

\section{Problem 2}

\subsection{Exact solution}

For the coaxial line,
\[
Z'=4+\jmath785.398163~\Omega/\mathrm{m},
\qquad
Y'=0.001+\jmath0.314159265~\mathrm{S/m}.
\]
The exact values are
\begin{align}
Z_0&=\sqrt{Z'/Y'}
=50.00017478-\jmath0.04774583~\Omega,\\
\gamma&=\sqrt{Z'Y'}
=0.064999970+\jmath15.707970430~\mathrm{m}^{-1}.
\end{align}
Therefore the exact attenuation is $0.064999970~\mathrm{Np/m}$ or
$0.56458974~\mathrm{dB/m}$.

\subsection{Low-loss model}

The two small parameters are
\[
\varepsilon_R=\frac{R}{\omega L}=0.005093,
\qquad
\varepsilon_G=\frac{G}{\omega C}=0.003183.
\]
Both are much less than one.  The first-order low-loss model yields
\begin{align}
Z_{0,\mathrm{LL}}&=\sqrt{L/C}=50.000~\Omega,\\
\alpha_{\mathrm{LL}}&=\frac{R}{2}\sqrt{\frac{C}{L}}
+\frac{G}{2}\sqrt{\frac{L}{C}}
=0.065000~\mathrm{Np/m},\\
\beta_{\mathrm{LL}}&=\omega\sqrt{LC}=15.70796327~\mathrm{rad/m}.
\end{align}
The comparison below uses magnitude for the exact characteristic impedance.

\begin{center}
\renewcommand{\arraystretch}{1.22}
\begin{tabular}{lrrr}
\toprule
 & Exact & Approximate & Approximation error \\
\midrule
$|Z_0|$ ($\Omega$) & 50.00019757 & 50.00000000 & $0.000395\%$ \\
$\alpha$ (Np/m) & 0.064999970 & 0.065000000 & $0.0000456\%$ \\
$\beta$ (rad/m) & 15.70797043 & 15.70796327 & $0.0000456\%$ \\
\bottomrule
\end{tabular}
\end{center}
The approximation is more than adequate at 500 MHz.

\subsection{Velocity and wavelength}

The exact phase values are
\begin{align}
v_p&=\frac{\omega}{\beta}=1.99999909\times10^8~\mathrm{m/s},\\
\lambda&=\frac{2\pi}{\beta}=0.399999818~\mathrm{m}.
\end{align}
The low-loss equations return $2.00000000\times10^8~\mathrm{m/s}$ and
$0.400000000~\mathrm{m}$, respectively.

\answer{$Z_0\approx50~\Omega$, $\alpha\approx0.065~\mathrm{Np/m}$,
$\beta\approx15.708~\mathrm{rad/m}$, $v_p\approx2.00\times10^8~\mathrm{m/s}$,
and $\lambda\approx0.400~\mathrm{m}$.}

\section{Problem 3}

\MatchedLineDiagram{v_s(t)=5\cos(2\pi\,2\!\times\!10^8t)}{75~\Omega}{50~\Omega}{3~\mathrm{m}}{R_L=50~\Omega}

\subsection{Forward-wave phasors}

Since $R_L=Z_0$, the reflection coefficient is zero:
\[
\Gamma_L=\frac{R_L-Z_0}{R_L+Z_0}=0.
\]
The line input also equals $50~\Omega$.  Thus
\begin{align*}
\phasor V(0)&=5\frac{50}{75+50}=2\angle0\degree~\mathrm{V},\\
\phasor I(0)&=\frac{\phasor V(0)}{50}=0.04\angle0\degree~\mathrm{A}.
\end{align*}
The wavelength is $v_p/f=1~\mathrm{m}$, giving $\beta=2\pi~\mathrm{rad/m}$.
Therefore
\begin{align}
\phasor V(z)&=2\ee^{-\jmath2\pi z}~\mathrm{V},\\
\phasor I(z)&=0.04\ee^{-\jmath2\pi z}~\mathrm{A}.
\end{align}
Returning to real time,
\begin{align}
v(z,t)&=2\cos(4\pi\times10^8t-2\pi z)~\mathrm{V},\\
i(z,t)&=0.04\cos(4\pi\times10^8t-2\pi z)~\mathrm{A}.
\end{align}

\subsection{Electrical length and load phase}

The normalized length and electrical length are
\[
\frac{\ell}{\lambda}=3,
\qquad \theta=\beta\ell=2\pi(3)=6\pi~\mathrm{rad}=1080\degree.
\]
Modulo $360\degree$, the load phase is $0\degree$.  Consequently,
\[
v_L(t)=2\cos(4\pi\times10^8t)~\mathrm{V},
\]
which is in phase with the ideal source waveform but smaller in amplitude
because of the $75~\Omega$--$50~\Omega$ voltage division.

\answer{The line contains exactly three wavelengths.  The load voltage is
$2~\mathrm{V}$ peak and has zero net phase shift relative to $v_s(t)$.}

\section{Problem 4}

\MatchedLineDiagram{v_s(t)=10\cos(2\pi\,5\!\times\!10^7t)}{25~\Omega}{100~\Omega}{6~\mathrm{m}}{R_L=100~\Omega}

\subsection{Matched-load fields}

The air-filled line has
\[
\lambda=\frac{3\times10^8}{5\times10^7}=6~\mathrm{m},
\qquad
\beta=\frac{2\pi}{6}=\frac{\pi}{3}~\mathrm{rad/m},
\qquad
\ell=\lambda.
\]
With $R_L=Z_0$, the input is $100~\Omega$ and
\[
V_0^+=10\frac{100}{25+100}=8~\mathrm{V},
\qquad I_0^+=0.08~\mathrm{A}.
\]
It follows that
\begin{align}
v(z,t)&=8\cos\left(\pi\times10^8t-\frac{\pi}{3}z\right)~\mathrm{V},\\
i(z,t)&=0.08\cos\left(\pi\times10^8t-\frac{\pi}{3}z\right)~\mathrm{A}.
\end{align}
The average power delivered to the matched load is
\begin{equation}
P_L=\frac{1}{2}\frac{|V_0^+|^2}{R_L}
=\frac{8^2}{2(100)}=0.320~\mathrm{W}.
\end{equation}

\subsection{Open-circuit case}

If the load is open, $\Gamma_L=+1$; a backward wave must be added.  Since
$\ell=\lambda$, the input repeats the open circuit and the source current is
zero.  Hence $\phasor V(0)=10~\mathrm{V}$.  Equal forward and backward wave
amplitudes of $5~\mathrm{V}$ satisfy this condition:
\begin{align}
\phasor V(z)&=5\ee^{-\jmath\beta z}+5\ee^{\jmath\beta z}
=10\cos(\beta z),\\
\phasor I(z)&=\frac{5}{100}\left(\ee^{-\jmath\beta z}
-\ee^{\jmath\beta z}\right)=-\jmath0.10\sin(\beta z).
\end{align}
Thus
\begin{align}
v(z,t)&=10\cos\left(\frac{\pi z}{3}\right)
\cos(\pi\times10^8t)~\mathrm{V},\\
i(z,t)&=0.10\sin\left(\frac{\pi z}{3}\right)
\sin(\pi\times10^8t)~\mathrm{A}.
\end{align}
The current vanishes at $z=0$ and $z=6~\mathrm{m}$, as expected for this
one-wavelength open-ended line.

\answer{The matched formulas give $P_L=0.320~\mathrm{W}$.  They do not remain
valid after opening the load: the reflected wave produces a standing-wave
voltage $10\cos(\pi z/3)\cos(\pi10^8t)$ V.}

\section*{Conclusion}
\addcontentsline{toc}{section}{Conclusion}
The exact lossy-line formulas were used in Problems 1 and 2, while the
reflection-free travelling-wave model was used in Problems 3 and 4.  The
low-loss coaxial approximation agrees with the exact calculation to much
better than $0.001\%$.  The open-circuit variant in Problem 4 demonstrates why
a load change must be accompanied by a reflection analysis.

\end{document}
